\section{Klasifikasi, Pemulihan, dan Pelaporan Kesalahan}
\label{sec:ch16-error}

\subsection{Klasifikasi Kesalahan}

\begin{table}[H]
\centering
\caption{Klasifikasi kesalahan dan fase pendeteksinya}
\label{tab:ch16-klasifikasi}
\begin{tabularx}{\textwidth}{@{}l l X@{}}
\toprule
\textbf{Kelas} & \textbf{Fase} & \textbf{Contoh} \\
\midrule
Leksikal & Lexer & karakter ilegal, string tidak ditutup \\
Sintaksis & Parser & titik koma hilang, kurung tidak seimbang \\
Semantik & Type checker & identifier tidak dideklarasikan, tipe salah \\
Validasi & Validator & \code{break} di luar loop, return hilang \\
\bottomrule
\end{tabularx}
\end{table}

Kesalahan \eng{run-time} seperti pembagian nol atau stack overflow
bukan kesalahan kompilasi; kesalahan itu muncul saat program
dijalankan. Guntoro membatasi jenis kesalahan program pada leksikal,
sintaksis, dan semantik.

\subsection{Contoh Kesalahan menurut Guntoro}

Contoh berikut diadaptasi ke subset bahasa C. Setiap contoh menunjukkan
jenis pelanggaran dan fase yang mendeteksinya.

\begin{table}[H]
\centering
\caption{Contoh kesalahan leksikal, sintaksis, dan semantik}
\label{tab:ch16-contoh-guntoro}
\begin{tabularx}{\textwidth}{@{}l l X@{}}
\toprule
\textbf{Jenis} & \textbf{Contoh} & \textbf{Penjelasan} \\
\midrule
Leksikal & \code{whle (a < b)} & kata kunci salah eja sehingga
menjadi identifier \\
Sintaksis & \code{a = x + (b * (c + d);} & kurung tidak seimbang \\
Semantik & \code{int siswa; siswa = "Yanuar";} & tipe string
diberikan ke variabel integer \\
Semantik & \code{b = b + 1;} & variabel \code{b} belum dideklarasikan \\
\bottomrule
\end{tabularx}
\end{table}

Pesan untuk setiap contoh harus menyebut baris, nama yang terlibat,
dan jenis pelanggaran, misalnya
\code{baris 3, kolom 12: tipe string tidak dapat ditetapkan ke siswa
bertipe int}.

\subsection{Pemulihan}

Setiap fase memiliki strategi pemulihan:
\begin{itemize}[leftmargin=*]
  \item lexer: lewati karakter ilegal;
  \item parser: sinkronisasi ke titik koma atau kurung tutup;
  \item type checker: gunakan \code{TYPE\_ERROR} dan lanjut;
  \item validator: catat kesalahan dan lanjut ke simpul berikutnya.
\end{itemize}

Pemulihan memungkinkan compiler melaporkan beberapa kesalahan dalam
satu proses. Namun, pemulihan yang agresif dapat menghasilkan pesan
palsu. Keseimbangan perlu diuji.

\begin{figure}[H]
\centering
\begin{tikzpicture}[node distance=6mm and 9mm]
  \node[decision, text width=2.6cm] (err) {Kesalahan};
  \node[phase, text width=2.6cm, below left=of err] (recover)
        {Pulihkan};
  \node[phase, text width=2.6cm, below right=of err] (stop)
        {Berhenti};
  \node[phase, text width=2.6cm, below=of recover] (continue)
        {Lanjutkan};
  \draw[arrow] (err) -- (recover);
  \draw[arrow] (err) -- (stop);
  \draw[arrow] (recover) -- (continue);
\end{tikzpicture}
\caption{Keputusan pemulihan setelah kesalahan}
\label{fig:ch16-recovery}
\end{figure}

\subsection{Format Diagnostik}

Format yang konsisten memudahkan pembacaan dan integrasi dengan
editor:
\[
\text{berkas}:\text{baris}:\text{kolom}: \text{error}: \text{pesan}.
\]
Contoh:
\begin{lstlisting}[language=bash]
tests/invalid/break_outside.sc:2:5: error: break hanya valid di dalam loop
\end{lstlisting}

\subsection{Batas Pelaporan}

Setelah sejumlah kesalahan, compiler sebaiknya berhenti dan
melaporkan ringkasan. Batas ini mencegah ledakan pesan dan membantu
pengguna fokus. Nilai batas dapat dikonfigurasi, misalnya 20
kesalahan.

\begin{table}[H]
\centering
\caption{Contoh pesan dari setiap fase}
\label{tab:ch16-pesan}
\begin{tabularx}{\textwidth}{@{}l X@{}}
\toprule
\textbf{Fase} & \textbf{Pesan} \\
\midrule
Lexer & \code{baris 1, kolom 1: string tidak ditutup} \\
Parser & \code{baris 2, kolom 10: diharapkan ';' setelah deklarasi} \\
Type checker & \code{baris 3, kolom 9: tidak dapat menetapkan float ke int} \\
Validator & \code{baris 2, kolom 5: break hanya valid di dalam loop} \\
\bottomrule
\end{tabularx}
\end{table}
